Agentic assistants for network and service operations are converging on the
same interface: the model is handed a set of tool (API) schemas and must pick
one and fill its arguments. In telecom this tool set is not small. A single
operator stack exposes order management, service and resource inventory,
catalogue, assurance, activation and billing interfaces, each with dozens of
operations. Published work has established that selection accuracy falls as the
tool set grows and that retrieving a small candidate set before prompting
recovers most of the loss~\cite{ragmcp,toollery}. It has also
established that a small model can internalise a \emph{fixed} tool catalogue by
fine-tuning, removing schemas from the prompt
entirely~\cite{qlora_tools}, and, in the opposite direction, that much of
what looks like a capability gap is really schema-name misalignment that can be
fixed by renaming tools rather than retraining them~\cite{patool}.

What is missing is a measurement of these choices on a tool catalogue that
telecom actually has. Existing telecom agent benchmarks evaluate
troubleshooting and dialogue over small, curated tool sets~\cite{telcoagent,ctbench,wirelessbench};
generic tool-routing studies use large but semantically scattered catalogues
where a nearest-neighbour lookup is usually enough. Standardised telecom
management APIs are different in a specific way: they are generated from a
common information model, so the same operation verb is applied to entities
whose names differ by one token. \texttt{listServiceCatalog},
\texttt{listResourceCatalog}, \texttt{listEntityCatalog} and
\texttt{listProductCatalog} are four different APIs, four different endpoints
and four different operational meanings, expressed in almost the same words.
We call such tools \emph{confusable}, and we show they are the majority case,
not a corner case.

To our knowledge this is the first tool-routing study on a standardised
telecom management catalogue at this scale. This work-in-progress paper makes
three contributions.
\textbf{(i)} We build a tool catalogue of 1,062 operations harvested from
79 published TM~Forum Open API specifications~\cite{tmforum}, and define a
confusability label that is computed from the specifications alone and is
independent of any retriever: 54.8\,\% of the catalogue is confusable.
\textbf{(ii)} We evaluate full-schema prompting, dense and lexical retrieval,
PA-Tool-style renaming, and a LoRA-tuned description-free 1.7B router on the
same queries, and report accuracy separately for confusable and distinct tools.
\textbf{(iii)} Because a wrong tool in this domain is not uniformly wrong, we
report a severity-weighted error cost and a destructive-misfire rate alongside
accuracy, and we score argument generation separately from selection so the two
failure modes do not mask each other.
